\section{Introduction}

Multi-objective planning problems arise when a decision maker must evaluate candidate plans using more than one objective. In many planning domains, a single scalar score is not sufficient to describe plan quality because objectives may conflict with one another. For example, one plan may minimize travel distance while another may reduce risk, energy consumption, or time. In this setting, comparing plans requires reasoning over objective vectors rather than individual values. Pareto dominance provides a principled way to perform this comparison by identifying plans that are not strictly worse than others across the selected objectives. The set of such nondominated plans forms the Pareto front and represents the trade-off boundary among competing objectives. 

This project builds on the Minimal Objectives Relevant to Nondominated Plans (MORNP) setting, where the goal is to reason about which objectives are necessary for preserving the dominance and nondominance structure among a collection of plans. In MORNP, plans can be separated into two structurally meaningful classes. The first class, denoted by $P^{+}$, contains nondominated plans that should remain relevant under the selected objective subset. The second class, denoted by $P^{-}$, contains dominated plans that should be excluded by the dominance relation. This distinction is important because the label of a plan is not simply an independent class label; it depends on how that plan relates to other plans in the objective space.

Generating large datasets for this type of reasoning is challenging. Standard synthetic data generation can create numerical vectors that appear realistic, but it may not preserve the dominance relationships that define the problem. A generated plan may have objective values that resemble the original data distribution while still violating the intended $P^{+}$ or $P^{-}$ structure. Therefore, a generative model for this domain must be evaluated not only by reconstruction quality or classification accuracy, but also by whether the generated samples remain valid under the dominance rules of the multi-objective planning problem.

To address this issue, this project proposes Gen-MOR, a generative modeling framework for learning and generating dominance-preserving multi-objective plans. Each plan is represented as an objective vector, and each vector is associated with a binary structural label indicating whether it belongs to $P^{+}$ or $P^{-}$. The central idea is to train neural generative models that can learn useful representations of these plan vectors, reconstruct existing samples, classify their dominance labels, and generate new synthetic plans that remain consistent with the MORNP structure.

The project evaluates three neural generative architectures: a plain autoencoder (AE), a variational autoencoder (VAE), and a conditional variational autoencoder (CVAE). The AE is used to study whether a compressed latent representation can preserve enough information to reconstruct objective vectors and separate $P^{+}$ from $P^{-}$. The VAE extends this idea by learning a regularized latent distribution from which new samples can be generated. The CVAE further incorporates the dominance label as a conditioning variable, allowing the model to generate samples targeted toward either the nondominated class or the dominated class. Together, these models provide a useful progression from representation learning to probabilistic generation and finally to class-controlled generation.

A key part of this project is the evaluation of generated samples using a dominance-aware metric. Standard metrics such as mean squared reconstruction error, accuracy, and F1 score are useful for measuring how well a model reconstructs or classifies existing data. However, they do not directly measure whether newly generated plans satisfy the relational structure required by Pareto dominance. For this reason, the project uses Dominance Validity Rate (DVR), a domain-specific metric that checks whether generated $P^{+}$ samples remain nondominated and whether generated $P^{-}$ samples remain dominated under the selected objective subset. DVR is therefore used as the main structural validity measure for the Gen-MOR framework.

The main contribution of this project is the formulation and experimental study of generative modeling for MORNP-style dominance-preserving plan generation. First, the project represents multi-objective plans as learnable objective vectors with dominance-based structural labels. Second, it trains AE, VAE and CVAE models to reconstruct, classify, and generate such plans. Third, it evaluates model behavior using both standard machine learning metrics and the dominance-specific DVR metric. Finally, it analyzes robustness across multiple random seeds to determine whether dominance-valid generation is stable or sensitive to training variation.

Overall, this work demonstrates that neural generative models can learn meaningful representations of multi-objective planning data, but it also shows that dominance-preserving generation is more difficult than ordinary reconstruction or classification. A model may achieve low reconstruction error while still generating samples that violate the intended dominance structure. This observation motivates the need for future generative models that incorporate dominance constraints directly into the learning objective, rather than treating dominance validity only as a post-hoc evaluation measure.